% Generated draft: requires mathematical and source review.
\documentclass[11pt]{article}
\usepackage[margin=25mm]{geometry}
\usepackage{fontspec}
\setmainfont{Latin Modern Roman}
\usepackage{amsmath,amssymb,mathtools,unicode-math}
\setmathfont{Latin Modern Math}
\usepackage{enumitem,needspace,xurl,hyperref}
\setlength{\parindent}{0pt}
\setlength{\parskip}{5pt}
\setlength{\emergencystretch}{3em}
\newcommand{\N}{\mathbb{N}}
\newcommand{\Z}{\mathbb{Z}}
\newcommand{\Q}{\mathbb{Q}}
\newcommand{\R}{\mathbb{R}}
\newcommand{\C}{\mathbb{C}}
\newcommand{\F}{\mathbb{F}}
\newcommand{\A}{\mathbb{A}}
\newcommand{\PP}{\mathbb{P}}
\newcommand{\E}{\mathbb{E}}
\newcommand{\eps}{\varepsilon}
\newcommand{\dd}{\,\mathrm{d}}
\begin{document}
\section*{UnsolvedMath ID: 30000391}
\textbf{Source problem code:} catalogue entry
\subsection{Definitions and mathematical statement}
Let \(d\geq1\) be an integer and let \(P\subseteq\R^d\) satisfy \(\operatorname{aff}(P)=\R^d\). No finiteness or countability assumption is imposed. A hyperplane spanned by \(P\) means an affine hyperplane generated by \(d\) affinely independent points of \(P\); explicitly, put
\[
H(P)=\{\operatorname{aff}(F):F\subseteq P,\ |F|=d,\ \dim\operatorname{aff}(F)=d-1\}.
\]
Hyperplanes are geometric sets, so different generating subsets defining the same hyperplane yield just one element of \(H(P)\).

The conjecture asks whether
\[
\forall d\geq1\ \forall P\subseteq\R^d\quad
\operatorname{aff}(P)=\R^d
\ \Longrightarrow\ 
\exists f:P\hookrightarrow H(P)\quad\forall p\in P,\ p\in f(p).
\]
Thus different points must be assigned different incident spanned hyperplanes, with no requirement that every hyperplane be used. Equivalently, in the bipartite incidence graph with parts \(P\) and \(H(P)\), is there a matching saturating \(P\)? The hyperplanes are affine, not necessarily linear through the origin, and not all incident hyperplanes are allowed unless they are spanned by \(P\).

The original short extraction omitted the full-affine-span hypothesis; the cleaned source-verified statement restores it, and the displayed conjecture retains it as substantive. The supplied literature triage checked 2026-08-21 treats the unrestricted infinite assertion as conjectural, noting that the cited finite Hall--matroid result does not prove it. Finite-case matching arguments are not silently applied to arbitrary infinite neighborhoods or uncountable point sets.
\subsection{Short English statement}
Can every affinely spanning point set in \(\R^d\), including arbitrary infinite sets, assign each point a different hyperplane spanned by the set and containing that point? This is an incidence matching that must cover all points.
\subsection{Sources}
\begin{itemize}
\item Local UnsolvedMath registry, permanent ID 30000391, catalogue entry, \emph{Infinite Transversals from Points to Spanned Hyperplanes}.
\item \emph{Combinatorics} (2006), Oberwolfach Reports, pp. 5--72, DOI \url{https://doi.org/10.4171/owr/2006/01}.
\item Supplied reference labeled \emph{Transversal matroids and Hall's theorem}: \url{https://msp.org/pjm/1972/41-3/pjm-v41-n3-p04-p.pdf}.
\item Supplied statement verification reports comparison with the official report; supplied literature triage checked 2026-08-21 distinguishes the finite cited theorem from the infinite question.
\end{itemize}
\end{document}
